% GENERATED FILE. Edit canonical lessons, then run npm run generate:books.
% canonical-source-hash: fnv1a64:3bc920142c75eff5
% canonical-lessons: FR-C479-blesser, FR-C479-coupure, FR-C479-egratignure, FR-C479-brulure, FR-C479-bosse

\chapter{Petits accidents}
\label{ch:fr-petits-accidents}
\hlchaptermodality{479}

\begin{chapteropening}
\textbf{By the end of this chapter:} \emph{I can talk about small injuries at home.}

Complete the last lesson of chapter 479: I can talk about small injuries at home.
\end{chapteropening}

\section[blesser]{blesser — to injure, to hurt; (se blesser) to hurt oneself}

\label{lesson:FR-C479-blesser}



\begin{quote}
Before the new one: say the French for nostril, then the French for temple.
\end{quote}

\subsection*{You'll want to know: blesser}
\textbf{blesser} — \textquotedblleft{}to injure, to hurt; (se blesser) to hurt oneself\textquotedblright{}.

\begin{quote}
Il s'est blessé au pied en jouant au foot.
\end{quote}

\begin{quote}
\emph{He hurt his foot playing football.}
\end{quote}

Often reflexive: \emph{se blesser à la main}. Also to hurt feelings: \emph{ses mots m'ont blessé}.

\begin{cousinweb}
\emph{blesser} comes from Old French \emph{blecier} \textquotedblleft{}to injure\textquotedblright{} (EtymDB), from Frankish \emph{blētjan} \textquotedblleft{}to bruise\textquotedblright{}, a Germanic word.
\end{cousinweb}

\begin{grammarlens}[title={What you've built}]
One more word for telling, explaining and describing at B1.
\end{grammarlens}

\subsection*{Your turn}
\begin{itemize}
\raggedright
  \item \emph{Say it:} \emph{blesser}
  \item \emph{Say it:} \emph{blesser}, once more
  \item \emph{Say it:} a sentence of your own with \emph{blesser}
  \item \emph{Recall:} say the French for nostril, then the French for temple, then say \emph{blesser} again
\end{itemize}

\begin{tcolorbox}[breakable,colback=teal!4,colframe=teal!35!black,title={Before you move on}]
What does blesser mean? (\textquotedblleft{}To injure, to hurt; (se blesser) to hurt oneself\textquotedblright{}.) Say it once more.
\end{tcolorbox}
\section[la coupure]{la coupure — cut}

\label{lesson:FR-C479-coupure}



\begin{quote}
Before the new one: say the French for temple, then the French for to injure.
\end{quote}

\subsection*{You'll want to know: la coupure}
\textbf{la coupure} — \textquotedblleft{}cut\textquotedblright{}.

\begin{quote}
J'ai une petite coupure au doigt.
\end{quote}

\begin{quote}
\emph{I have a small cut on my finger.}
\end{quote}

Also a break in a supply: \emph{une coupure de courant} — a power cut.

\begin{cousinweb}
\emph{la coupure} is built on \emph{couper} \textquotedblleft{}to cut\textquotedblright{} with \emph{-ure} (EtymDB); \emph{couper} comes from \emph{coup} \textquotedblleft{}blow\textquotedblright{}, from Latin \emph{colaphus} \textquotedblleft{}a punch\textquotedblright{}.
\end{cousinweb}

\begin{grammarlens}[title={What you've built}]
One more word for telling, explaining and describing at B1.
\end{grammarlens}

\subsection*{Your turn}
\begin{itemize}
\raggedright
  \item \emph{Say it:} \emph{la coupure}
  \item \emph{Say it:} \emph{la coupure}, once more
  \item \emph{Say it:} a sentence of your own with \emph{la coupure}
  \item \emph{Recall:} say the French for temple, then the French for to injure, then say \emph{la coupure} again
\end{itemize}

\begin{tcolorbox}[breakable,colback=teal!4,colframe=teal!35!black,title={Before you move on}]
What does la coupure mean? (\textquotedblleft{}Cut\textquotedblright{}.) Say it once more.
\end{tcolorbox}
\section[l'égratignure]{l'égratignure — scratch, graze}

\label{lesson:FR-C479-egratignure}



\begin{quote}
Before the new one: say the French for to injure, then the French for cut.
\end{quote}

\subsection*{You'll want to know: l'égratignure}
\textbf{l'égratignure} — \textquotedblleft{}scratch, graze\textquotedblright{}.

\begin{quote}
Le chat m'a fait une égratignure sur le bras.
\end{quote}

\begin{quote}
\emph{The cat gave me a scratch on my arm.}
\end{quote}

Feminine. Lighter than \emph{une coupure}.

\begin{cousinweb}
\emph{l'égratignure} is built on \emph{égratigner} \textquotedblleft{}to scratch\textquotedblright{}, itself built on \emph{gratter} \textquotedblleft{}to scratch\textquotedblright{}, which comes from a Germanic word, Frankish \emph{krattōn}.
\end{cousinweb}

\begin{grammarlens}[title={What you've built}]
One more word for telling, explaining and describing at B1.
\end{grammarlens}

\subsection*{Your turn}
\begin{itemize}
\raggedright
  \item \emph{Say it:} \emph{l'égratignure}
  \item \emph{Say it:} \emph{l'égratignure}, once more
  \item \emph{Say it:} a sentence of your own with \emph{l'égratignure}
  \item \emph{Recall:} say the French for to injure, then the French for cut, then say \emph{l'égratignure} again
\end{itemize}

\begin{tcolorbox}[breakable,colback=teal!4,colframe=teal!35!black,title={Before you move on}]
What does l'égratignure mean? (\textquotedblleft{}Scratch, graze\textquotedblright{}.) Say it once more.
\end{tcolorbox}
\section[la brûlure]{la brûlure — burn}

\label{lesson:FR-C479-brulure}



\begin{quote}
Before the new one: say the French for cut, then the French for scratch.
\end{quote}

\subsection*{You'll want to know: la brûlure}
\textbf{la brûlure} — \textquotedblleft{}burn\textquotedblright{}.

\begin{quote}
Elle s'est fait une brûlure en sortant le gâteau du four.
\end{quote}

\begin{quote}
\emph{She burnt herself taking the cake out of the oven.}
\end{quote}

Built on \emph{brûler} \textquotedblleft{}to burn\textquotedblright{}.

\begin{cousinweb}
\emph{la brûlure} is built on \emph{brûler} \textquotedblleft{}to burn\textquotedblright{} with \emph{-ure}. The origin of \emph{brûler} is not settled: Latin \emph{ustulāre} \textquotedblleft{}to scorch\textquotedblright{} and a Germanic verb for burning have both been proposed.
\end{cousinweb}

\begin{grammarlens}[title={What you've built}]
One more word for telling, explaining and describing at B1.
\end{grammarlens}

\subsection*{Your turn}
\begin{itemize}
\raggedright
  \item \emph{Say it:} \emph{la brûlure}
  \item \emph{Say it:} \emph{la brûlure}, once more
  \item \emph{Say it:} a sentence of your own with \emph{la brûlure}
  \item \emph{Recall:} say the French for cut, then the French for scratch, then say \emph{la brûlure} again
\end{itemize}

\begin{tcolorbox}[breakable,colback=teal!4,colframe=teal!35!black,title={Before you move on}]
What does la brûlure mean? (\textquotedblleft{}Burn\textquotedblright{}.) Say it once more.
\end{tcolorbox}
\section[la bosse]{la bosse — bump, lump}

\label{lesson:FR-C479-bosse}



\begin{quote}
Before the new one: say the French for scratch, then the French for burn.
\end{quote}

\subsection*{You'll want to know: la bosse}
\textbf{la bosse} — \textquotedblleft{}bump, lump\textquotedblright{}.

\begin{quote}
Après sa chute, l'enfant a une bosse sur le front.
\end{quote}

\begin{quote}
\emph{After his fall, the child has a bump on his forehead.}
\end{quote}

Also a camel's hump: \emph{la bosse du chameau}.

\begin{cousinweb}
\emph{la bosse} goes back to Old French \emph{boce} (EtymDB). Its earlier origin is not settled; a Vulgar Latin \emph{bottia} \textquotedblleft{}swelling\textquotedblright{} is sometimes proposed.
\end{cousinweb}

\begin{grammarlens}[title={What you've built}]
One more word for telling, explaining and describing at B1.
\end{grammarlens}

\subsection*{Your turn}
\begin{itemize}
\raggedright
  \item \emph{Say it:} \emph{la bosse}
  \item \emph{Say it:} \emph{la bosse}, once more
  \item \emph{Say it:} a sentence of your own with \emph{la bosse}
  \item \emph{Recall:} say the French for scratch, then the French for burn, then say \emph{la bosse} again
\end{itemize}

\begin{tcolorbox}[breakable,colback=teal!4,colframe=teal!35!black,title={Before you move on}]
What does la bosse mean? (\textquotedblleft{}Bump, lump\textquotedblright{}.) Say it once more.
\end{tcolorbox}
